% GENERATED FILE. Edit canonical lessons, then run npm run generate:books.
% canonical-source-hash: fnv1a64:a1dbb5b73cd12552
% canonical-lessons: IT-C430-piantare, IT-C430-sacco-a-pelo, IT-C430-roulotte, IT-C430-camper, IT-C430-zanzariera

\chapter{In campeggio}
\label{ch:it-in-campeggio}
\hlchaptermodality{430}

\begin{chapteropening}
\textbf{By the end of this chapter:} \emph{I can talk about camping: pitching a tent, the sleeping bag, the caravan, the camper van and the mosquito net.}

Complete the last lesson of chapter 430: I can talk about camping: pitching a tent, the sleeping bag, the caravan, the camper van and the mosquito net.
\end{chapteropening}

\section[piantare]{piantare — to pitch (a tent); to plant}

\label{lesson:IT-C430-piantare}



\begin{quote}
Before the new one: say the Italian for the guesthouse, then the Italian for the campsite.
\end{quote}

\subsection*{You'll want to know: piantare}
\textbf{piantare} — \textquotedblleft{}to pitch (a tent); to plant\textquotedblright{}.

\begin{quote}
Piantiamo la tenda vicino al fiume.
\end{quote}

\begin{quote}
\emph{We pitch the tent near the river.}
\end{quote}

Also \emph{piantare un albero} — to plant a tree; in speech, \emph{piantare qualcuno} is to leave someone.

\begin{cousinweb}
\emph{piantare} is inherited from Latin \emph{plantāre} 'to plant' (EtymDB), from \emph{planta} 'sole of the foot, young shoot'.
\end{cousinweb}

\begin{grammarlens}[title={What you've built}]
One more word for telling, explaining and describing at B1.
\end{grammarlens}

\subsection*{Your turn}
\begin{itemize}
\raggedright
  \item \emph{Say it:} \emph{piantare}
  \item \emph{Say it:} \emph{piantare}, once more
  \item \emph{Say it:} a sentence of your own with \emph{piantare}
  \item \emph{Recall:} say the Italian for the guesthouse, then the Italian for the campsite, then say \emph{piantare} again
\end{itemize}

\begin{tcolorbox}[breakable,colback=teal!4,colframe=teal!35!black,title={Before you move on}]
What does piantare mean? (\textquotedblleft{}To pitch (a tent); to plant\textquotedblright{}.) Say it once more.
\end{tcolorbox}
\section[il sacco a pelo]{il sacco a pelo — the sleeping bag}

\label{lesson:IT-C430-sacco-a-pelo}



\begin{quote}
Before the new one: say the Italian for the campsite, then the Italian for to pitch.
\end{quote}

\subsection*{You'll want to know: il sacco a pelo}
\textbf{il sacco a pelo} — \textquotedblleft{}the sleeping bag\textquotedblright{}.

\begin{quote}
Di notte fa freddo, porta il sacco a pelo.
\end{quote}

\begin{quote}
\emph{It is cold at night, bring the sleeping bag.}
\end{quote}

Literally a sack with fur: the first ones were lined with it.

\begin{cousinweb}
\emph{sacco a pelo} is built from \emph{sacco} 'sack', inherited from Latin \emph{saccus}, from Greek \emph{sákkos}, and \emph{pelo} 'hair, fur', inherited from Latin \emph{pilus} (EtymDB).
\end{cousinweb}

\begin{grammarlens}[title={What you've built}]
One more word for telling, explaining and describing at B1.
\end{grammarlens}

\subsection*{Your turn}
\begin{itemize}
\raggedright
  \item \emph{Say it:} \emph{il sacco a pelo}
  \item \emph{Say it:} \emph{il sacco a pelo}, once more
  \item \emph{Say it:} a sentence of your own with \emph{il sacco a pelo}
  \item \emph{Recall:} say the Italian for the campsite, then the Italian for to pitch, then say \emph{il sacco a pelo} again
\end{itemize}

\begin{tcolorbox}[breakable,colback=teal!4,colframe=teal!35!black,title={Before you move on}]
What does il sacco a pelo mean? (\textquotedblleft{}The sleeping bag\textquotedblright{}.) Say it once more.
\end{tcolorbox}
\section[la roulotte]{la roulotte — the caravan}

\label{lesson:IT-C430-roulotte}



\begin{quote}
Before the new one: say the Italian for to pitch, then the Italian for the sleeping bag.
\end{quote}

\subsection*{You'll want to know: la roulotte}
\textbf{la roulotte} — \textquotedblleft{}the caravan\textquotedblright{}.

\begin{quote}
I nonni passano l'estate in roulotte.
\end{quote}

\begin{quote}
\emph{My grandparents spend the summer in a caravan.}
\end{quote}

Said the French way; towed behind a car. Plural unchanged.

\begin{cousinweb}
\emph{roulotte} was borrowed from French \emph{roulotte} (EtymDB), built on \emph{rouler} 'to roll', which goes back to Latin \emph{rota} 'wheel'.
\end{cousinweb}

\begin{grammarlens}[title={What you've built}]
One more word for telling, explaining and describing at B1.
\end{grammarlens}

\subsection*{Your turn}
\begin{itemize}
\raggedright
  \item \emph{Say it:} \emph{la roulotte}
  \item \emph{Say it:} \emph{la roulotte}, once more
  \item \emph{Say it:} a sentence of your own with \emph{la roulotte}
  \item \emph{Recall:} say the Italian for to pitch, then the Italian for the sleeping bag, then say \emph{la roulotte} again
\end{itemize}

\begin{tcolorbox}[breakable,colback=teal!4,colframe=teal!35!black,title={Before you move on}]
What does la roulotte mean? (\textquotedblleft{}The caravan\textquotedblright{}.) Say it once more.
\end{tcolorbox}
\section[il camper]{il camper — the camper van, the motorhome}

\label{lesson:IT-C430-camper}



\begin{quote}
Before the new one: say the Italian for the sleeping bag, then the Italian for the caravan.
\end{quote}

\subsection*{You'll want to know: il camper}
\textbf{il camper} — \textquotedblleft{}the camper van, the motorhome\textquotedblright{}.

\begin{quote}
Con il camper dormiamo dove vogliamo.
\end{quote}

\begin{quote}
\emph{With the camper van we sleep where we like.}
\end{quote}

\emph{Il camperista} is the owner. Plural unchanged.

\begin{cousinweb}
\emph{camper} was borrowed from English \emph{camper} (EtymDB), built on \emph{camp}, which goes back to Latin \emph{campus} 'field'.
\end{cousinweb}

\begin{grammarlens}[title={What you've built}]
One more word for telling, explaining and describing at B1.
\end{grammarlens}

\subsection*{Your turn}
\begin{itemize}
\raggedright
  \item \emph{Say it:} \emph{il camper}
  \item \emph{Say it:} \emph{il camper}, once more
  \item \emph{Say it:} a sentence of your own with \emph{il camper}
  \item \emph{Recall:} say the Italian for the sleeping bag, then the Italian for the caravan, then say \emph{il camper} again
\end{itemize}

\begin{tcolorbox}[breakable,colback=teal!4,colframe=teal!35!black,title={Before you move on}]
What does il camper mean? (\textquotedblleft{}The camper van, the motorhome\textquotedblright{}.) Say it once more.
\end{tcolorbox}
\section[la zanzariera]{la zanzariera — the mosquito net, the fly screen}

\label{lesson:IT-C430-zanzariera}



\begin{quote}
Before the new one: say the Italian for the caravan, then the Italian for the camper van.
\end{quote}

\subsection*{You'll want to know: la zanzariera}
\textbf{la zanzariera} — \textquotedblleft{}the mosquito net, the fly screen\textquotedblright{}.

\begin{quote}
Chiudi la zanzariera, entrano gli insetti.
\end{quote}

\begin{quote}
\emph{Close the fly screen, the insects are getting in.}
\end{quote}

Built on \emph{la zanzara}, the mosquito.

\begin{cousinweb}
\emph{zanzariera} is built on \emph{zanzara} 'mosquito', which is inherited from Late Latin \emph{zinzala} (EtymDB), a word that imitates the buzzing.
\end{cousinweb}

\begin{grammarlens}[title={What you've built}]
One more word for telling, explaining and describing at B1.
\end{grammarlens}

\subsection*{Your turn}
\begin{itemize}
\raggedright
  \item \emph{Say it:} \emph{la zanzariera}
  \item \emph{Say it:} \emph{la zanzariera}, once more
  \item \emph{Say it:} a sentence of your own with \emph{la zanzariera}
  \item \emph{Recall:} say the Italian for the caravan, then the Italian for the camper van, then say \emph{la zanzariera} again
\end{itemize}

\begin{tcolorbox}[breakable,colback=teal!4,colframe=teal!35!black,title={Before you move on}]
What does la zanzariera mean? (\textquotedblleft{}The mosquito net, the fly screen\textquotedblright{}.) Say it once more.
\end{tcolorbox}
